\section{Diagnostic examples and common traps}

Conditional probability is conceptually simple once the reference space is clear,
but it is easy to make mistakes when the roles of events are not stated
explicitly. This section collects common traps and shows how to diagnose them.

The guiding question is always:
\[
\text{what is the information, and what is being measured inside that information?}
\]

\subsection*{Trap 1: confusing \(P(A\mid B)\) with \(P(B\mid A)\)}

The conditional probability
\[
P(A\mid B)
\]
means that \(B\) is the information and \(A\) is the event measured inside that
information. The conditional probability
\[
P(B\mid A)
\]
reverses these roles.

These two numbers are usually different.

\begin{examplebox}
Let
\[
A=\{\text{a person is a mathematics student}\},
\]
and
\[
B=\{\text{a person takes a probability course}\}.
\]
The probability \(P(B\mid A)\) asks: among mathematics students, what proportion
take a probability course? The probability \(P(A\mid B)\) asks: among students in
the probability course, what proportion are mathematics students?

These are different reference groups, so the answers need not be the same.
\end{examplebox}

Diagnostic question:
\[
\text{Which event is after the vertical bar?}
\]
That event is the new reference space.

\subsection*{Trap 2: ignoring base rates}

Bayesian problems often involve a rare event and an imperfect observation. The
observation may be strong evidence, but the prior probability still matters.
Ignoring the prior probability is called ignoring the base rate.

Suppose a condition occurs in \(1\%\) of a population. A test is positive for
\(99\%\) of people with the condition and for \(5\%\) of people without the
condition. A positive test does not mean the person has probability \(99\%\) of
having the condition. The number \(99\%\) is
\[
P(\text{positive}\mid\text{condition}).
\]
The desired probability after a positive test is
\[
P(\text{condition}\mid\text{positive}).
\]
These are different.

Diagnostic question:
\[
\text{How common was the case before the evidence was observed?}
\]
Bayes' theorem combines the prior probability with the likelihood of the
observation.

\subsection*{Trap 3: treating disjoint events as independent}

If two events cannot occur together, they are disjoint:
\[
A\cap B=\varnothing.
\]
If both have positive probability, then learning that one occurred makes the
other impossible. Therefore they are dependent, not independent.

For example, in one die roll, let
\[
A=\{\text{the result is even}\}=\{2,4,6\},
\]
and
\[
B=\{\text{the result is odd}\}=\{1,3,5\}.
\]
The events are disjoint. If we know \(B\), then the probability of \(A\) becomes
zero:
\[
P(A\mid B)=0.
\]
But before knowing \(B\),
\[
P(A)=\frac12.
\]
The information changes the probability, so the events are dependent.

Diagnostic question:
\[
\text{If one event occurs, does the other become impossible?}
\]
If yes, then positive-probability events are not independent.

\subsection*{Trap 4: assuming independence without justification}

Independence is not created by notation. Writing a sample space as a product,
such as
\[
\{H,T\}^2
\]
or
\[
\{1,2,3,4,5,6\}^2,
\]
describes possible ordered outcomes. It does not by itself prove that the stages
are independent.

Independence is a property of the probability assignment. It says that knowing
one event occurred does not change the probability of another.

\begin{examplebox}
A mechanism produces two symbols, each either \(H\) or \(T\). The possible
outputs are
\[
HH,HT,TH,TT.
\]
If the mechanism chooses the first symbol fairly and then always repeats it, the
only outputs with positive probability are \(HH\) and \(TT\). The two positions
are not independent, even though the notation has two coordinates.
\end{examplebox}

Diagnostic question:
\[
\text{Where in the model is independence stated or justified?}
\]
If it is not stated or implied by the mechanism, do not use product rules as if
it were automatic.

\subsection*{Trap 5: using total probability without a partition}

The law of total probability requires cases that are mutually exclusive and
exhaustive. If the cases overlap, adding contributions may double-count. If the
cases do not cover the whole sample space, the sum may miss part of the event.

Before writing
\[
P(A)=\sum_i P(A\mid B_i)P(B_i),
\]
check that the events \(B_i\) form a partition.

Diagnostic questions:
\begin{itemize}
    \item Can two of the cases occur at the same time?
    \item Is every possible outcome covered by one of the cases?
    \item Do all conditioning events have positive probability?
\end{itemize}

\subsection*{Trap 6: forgetting the denominator in Bayes' theorem}

In Bayes' theorem, the denominator is the total probability of the observation.
It includes all ways the observation can occur. Omitting some cases changes the
reference space and gives the wrong posterior probability.

For a partition \(B_1,
\ldots,B_n\), the denominator is
\[
P(A)=\sum_i P(A\mid B_i)P(B_i).
\]
It is not merely the likelihood under the case we are interested in.

Diagnostic question:
\[
\text{What are all possible sources of the observation?}
\]
The denominator must include them all.

\subsection*{Trap 7: treating a conditional probability as a causal statement}

If
\[
P(A\mid B)>P(A),
\]
then \(B\) is informative about \(A\). But this alone does not prove that \(B\)
causes \(A\). Conditional probability describes how probabilities change under
information. Causal interpretation requires additional structure.

For example, if students who attend review sessions pass more often, then
attendance is associated with passing. But from this probability statement alone
we cannot conclude that attendance caused the higher pass rate. Students who
attend may differ in other ways.

Diagnostic question:
\[
\text{Am I making a probability statement or a causal claim?}
\]
These are not the same kind of statement.

\subsection*{Trap 8: conditioning on an event of probability zero}

The elementary formula
\[
P(A\mid B)=\frac{P(A\cap B)}{P(B)}
\]
requires
\[
P(B)>0.
\]
If \(P(B)=0\), the denominator vanishes and the formula is not defined.

In finite examples this often corresponds to impossible information. For a die
roll, conditioning on the event ``the result is \(8\)'' is not meaningful in the
usual die model because that event is empty.

Diagnostic question:
\[
\text{Does the conditioning event have positive probability in this model?}
\]

\subsection*{A diagnostic workflow}

When a conditional probability problem appears, use the following workflow.

\begin{enumerate}
    \item Identify the original sample space or probability model.
    \item Identify the information event: this is the event after the vertical
    bar.
    \item Identify the target event: this is the event before the vertical bar.
    \item Check that the information event has positive probability.
    \item Replace the active space of reasoning by the information event.
    \item Locate the target event inside that information event by taking the
    intersection.
    \item Decide whether the problem needs direct conditioning, independence,
    total probability, or Bayes' theorem.
\end{enumerate}

\begin{theorembox}
Most mistakes with conditional probability are mistakes about reference spaces.
Before calculating, say explicitly: ``I am reasoning inside this event.''
\end{theorembox}
